\subsection{Comparison across families}
\label{sec:comparison}

\begin{table*}[!tb]
\centering\footnotesize\setlength\tabcolsep{4pt}\renewcommand{\arraystretch}{1.12}
\caption{What each method family shows on its own, what it does not, the tests that license a stronger
reading (none a prerequisite; Table~\ref{tab:vcov}), and where this corpus stops, with $k/n$ over the family's own records (the cross-family
counts are Table~\ref{tab:vcov}). Geometry is task-linked geometry only (\S\ref{sec:models}); the sub-types of decoding and intervention are discussed in \S\ref{sec:decoding} and \S\ref{sec:intervention}.}
\label{tab:families}
\begin{tabular}{@{}L{2.3cm}L{3.4cm}L{2.8cm}L{3.25cm}L{3.25cm}@{}}
\toprule
Family & Shows on its own & Does not show & Stronger reading needs & In this corpus \\
\midrule
\textbf{Structural and geometric} (task-linked) & the shape of the space: gap, alignment, effective rank, similarity & what any direction means; no concept & comparison model; random-initialised or shuffled baseline & comparison model \Vcomp{}/\Vcompn{}; calibration \Vncalgeo{}/\Vncalgeon{} \\
\addlinespace[3pt]
\textbf{Supervised concept mapping and decodability} & a concept can be read out of $z_\lambda$ by a decoder (Eq.~\ref{eq:probe}) & that the model uses it & permuted labels; random encoder; control task; fixed decoder class & permutation null \Vpermdec{}/\Vpermdecn{}; random encoder \Vlrndec{}/\Vlrndecn{}; selectivity not yet measurable \\
\addlinespace[3pt]
\textbf{Feature decomposition and discovery} & a unit's activation tracks the named content (Eq.~\ref{eq:sae}) & that the unit is stable, complete or used & validated names; cross-seed matching; reconstruction & semantic validation \Vsem{}/\Vsemn{}; cross-seed identity none \\
\addlinespace[3pt]
\textbf{Attribution, localisation and tracing} & where information is decodable or attended ($\ell$, $p$, $h$) & that the location is used & occlusion or patching confirmation; random-component comparison & input corroboration \Vinp{}/\Vinpn{}; component perturbation \Vinttr{}/\Vinttrn{} \\
\addlinespace[3pt]
\textbf{Representation intervention and control} & the output is sensitive to $z_\lambda$ (Eq.~\ref{eq:effect}) & that the named concept, not a correlated direction, caused the change & matched control directions; equal-size random units; sham transformations & matched control \Vspca{}--\Vspc{}/\Vspcn{}; no dose--response test \\
\midrule
\textbf{Intrinsic and concept-aligned designs} (tag) & the label depends on $z_\lambda$ only through the concept layer & that nothing else leaks through & leakage test; random concept projections; concept-layer intervention & concept-layer control \Vintrany{}/\Vintrn{}; leakage rarely measured \\
\bottomrule
\end{tabular}
\end{table*}
\begin{table*}[!tb]
\centering\footnotesize\setlength\tabcolsep{2.5pt}\renewcommand{\arraystretch}{1.05}
\caption{What each family costs and how far this corpus has taken it. \emph{Concept annotation}: whether
the design needs a concept vocabulary, and how much labelling. \emph{Loci}: where the family has been
applied to medical models (parentheses: reached only in general-domain work). \emph{Compute}: beyond a
forward pass. \emph{Readable}: whether a clinician can read the artefact directly. \emph{Base.}: records
comparing a medical with a general model on the same data. \emph{Multi.}: more than one model or dataset.
The structural row ($\dagger$) covers task-linked geometry only (\S\ref{sec:models}) and is not comparable
with the others; the \Nintrinsic{} intrinsic-design records are counted in their analytic family and
listed again in the tag row.}
\label{tab:famreq}
\begin{tabular}{@{}L{2.2cm}L{2.1cm}L{2.0cm}L{1.6cm}L{1.8cm}C{0.7cm}C{0.8cm}C{0.8cm}L{2.8cm}@{}}
\toprule
Family & Concept annotation & Loci (medical) & Compute & Readable & 1st yr & Base. & Multi. & Main failure mode \\
\midrule
Structural and geometric (task-linked)$^{\dagger}$ & none & Enc, Txt, Joint, Dec, (Conn) & low & no & 2022 & 11 & 29 & geometric quantity read as clinical meaning \\
Supervised concept mapping and decodability & required; high (labels) & Enc, Joint & low & partly (scores yes, probes no) & 2022 & 31 & 72 & decodability read as use \\
Feature decomposition and discovery & none or weak (post-hoc naming) & Enc, Dec, (Conn) & high (activations at scale, training) & yes, once named and validated & 2024 & 2 & 14 & unstable features and unvalidated names \\
Attribution, localisation and tracing & optional; low--moderate & Enc, Joint, Dec, Path, (Pre) & moderate & partly (maps, often unfaithful) & 2022 & 2 & 8 & localisation read as causal confirmation \\
Representation intervention and control & required for concept-\hspace{0pt}specific claims; moderate (exemplars) & Enc, Joint, Dec, (Txt, Conn) & mod.--high (repeated generation) & indirectly (behaviour change) & 2023 & 5 & 18 & uncontrolled, input-\hspace{0pt}dependent effects \\
\midrule
Intrinsic and concept-aligned designs (tag) & required; high & Enc (concept layer), Joint & low--moderate & yes & 2023 & 7 & 13 & leakage; concept invalidity \\
\bottomrule
\end{tabular}
\end{table*}

Table~\ref{tab:families} sets out what each family reports and what its artefact licenses;
Table~\ref{tab:famreq} gives what each costs and how far this corpus has taken it. Decoding is the
oldest and by far the largest family; discovery is two years old and intervention three. Two asymmetries
in Table~\ref{tab:famreq} are easily missed. Intervention gives the claims of greatest practical interest
at the highest cost and with the fewest controlled results, while the cheapest family addresses
availability, which is often mistaken for use. Readability and validity also pull apart: discovery and
intrinsic designs return the artefacts a clinician can read, yet they rest on an unvalidated naming or
leakage step.

\textbf{Combining families.} Three pipelines recur. \emph{Discover, then intervene} uses residual-stream
dictionary directions for steering or
patching~\cite{bouzid2025insights,nooralahzadeh2026universal,sadanandan2026mechanistically}, and carries
the caveats of \S\ref{sec:discovery} into the intervention. \emph{Trace, then intervene} exploits an
observational localisation through a training-time change, so far only in a general CLIP text tower on
radiology captions~\cite{abbasi2026layerspecific}. \emph{Decode, then erase, residualise or ablate}
probes for availability and then tests whether a downstream prediction depends on the
concept~\cite{kim2026encoding,nooralahzadeh2026sparse,weber2023posthoc}; it answers the shortcut
question of \S\ref{sec:applications} most directly, and its two core instances that test a downstream
dependence point in opposite directions, residualisation leaving a refitted finding probe
unchanged~\cite{kim2026encoding}$^{\dagger}$ while channel knockout collapses a finding's
score~\cite{nooralahzadeh2026sparse}. Reading one family's result as an answer to another family's
question is the commonest inferential error in this literature, and the evidence-type tag makes it
visible.

